\section{Outward propagation and an electron deficit}
\label{a:sec:propagation}

The conditional comparisons let us extend a prescribed electron source
from the scale $Z^{-1/3}$ to a small radius independent of $Z$.  The
result applies to ground states whose energy has a bounded offset from
the minimum over all particle numbers.

\begin{proposition}[Deficit inside a fixed small ball]
\label{a:prop:prop-deficit}
There is a universal constant $c_{\mathrm{def}}>0$ with the following
property.  For every $D_0\in[0,\infty)$ there is $s_*(D_0)>0$ such that,
for each $0<s\le s_*(D_0)$, there is an integer $Z_*(D_0,s)$ for which
every integer $Z\ge Z_*(D_0,s)$ and every normalized sector ground state
$\Psi$ satisfying
\[
 n\le3Z,\qquad q_n[\Psi]=E_n\le E+D_0
\]
obey
\begin{equation}
\label{a:eq:prop-deficit}
 Z-\E_\Psi\#\{i:|x_i|<s/4\}
 \ge c_{\mathrm{def}}s^{-3}>1.
\end{equation}
All thresholds are independent of the sector and of the choice of its
ground state.  In particular, admissible fixed radii $s$ can be taken
arbitrarily small.
\end{proposition}


To prove the deficit, we construct fields $u_j$ satisfying a lower bound
for their Laplacian. Alongside the nuclear term $-4\pi\nu$, this bound
uses the conditional density term $4\pi\mu_j$ inside radius $r_j$ and
the model reaction $4\pi k(u_j)_+^{3/2}$ outside. Each outward step
replaces that reaction by the actual density on one more annulus, then
forgets the current observation array. An error density records the
source omitted on exceptional data. A positive lower barrier and a small
expected error will force a positive residual charge at the terminal
scale. The construction also gives the following result at every
intermediate stopping index.

\begin{proposition}[Intermediate atomic subsolutions]
\label{prop:atomic-intermediate}
There are universal constants $\eps,B,C_{\rm inv}>0$ and $L_1>2$
with the following property. Fix $D_0<\infty$, then take $s>0$
sufficiently small and integer $Z$ sufficiently large as in
Proposition~\ref{a:prop:prop-deficit}. For every sector ground state in
that proposition use the common observations and master packets
\eqref{a:eq:cond-observations}--\eqref{a:eq:cond-master}, with
$r_0=\eps Z^{-1/3}$ and $s/2<r_m\le s$.
For each $0\le j\le m$ there exist jointly measurable
$\mathcal F_j$-measurable fields $u_j$ and densities $p_j^{\rm err}$
with the following properties. Put $d=|y|$, $r=r_j$, and
$f(t)=4\pi k(t_+)^{3/2}$, and define
\[
 b_r(y)=Bd^{-4}\left(1-\frac r{8d}\right)\quad(d\ge r).
\]
\begin{equation}
\label{a:eq:prop-invariant}
 \begin{split}
 \Delta u_j&\ge-4\pi\nu+4\pi\1_{\{d<r\}}\mu_j
       -4\pi p_j^{\mathrm{err}}+\1_{\{d\ge r\}}f(u_j),\\
 b_r\le u_j&\le C_{\mathrm{inv}}d^{-4}\quad(d\ge r),
 \qquad u_j=b_r\quad(d>L_1r).
 \end{split}
\end{equation}
The differential inequality is global in distributions.
The offsets $u_j-V$ are continuous and locally Lipschitz on all of
$\R^3$; $p_j^{\rm err}\ge0$ is bounded and supported in $|y|\le r_j$.
For each fixed system and step, the offsets have deterministic size
and Lipschitz bounds on every compact set; the errors have deterministic
size bounds. Universally,
\begin{equation}\label{eq:atomic-prefix-error}
 \E\int p_j^{\rm err}
 \le C r_0^{32}+C\sum_{i<j}r_i^{19/2-3w}
 \le C\left(r_0^{32}+r_j^{19/2-3w}\right),\qquad w=10^{-5}.
\end{equation}
In particular this is at most $C(s^{32}+s^{19/2-3w})$.
The expectations cannot in general be omitted for $j<m$; the terminal
fields are deterministic. The same master kernels are used at every
intermediate stopping index. The constants $\eps,B,C_{\rm inv},L_1$
are independent of any additional inverse-comparison parameters.
Such a comparison uses these same observations whenever its own
smallness condition on $s$ is met.
\end{proposition}

\begin{proof}[Proof of Propositions~\ref{prop:atomic-intermediate} and~\ref{a:prop:prop-deficit}]
Use the observations and master packets
\eqref{a:eq:cond-observations}--\eqref{a:eq:cond-master}, with $\eps$
fixed during initialization below and $Z$ large enough that $m\ge1$.
The resulting densities and fields satisfy
\begin{equation}
\label{a:eq:prop-master-recall}
 \E[\mu_j\mid\mathcal F_{j+1}]=\mu_{j+1},\qquad
 \E[H_j\mid\mathcal F_{j+1}]=H_{j+1}.
\end{equation}

Write $\kappa=1/8$. The barrier $b_r$ in the proposition is positive
on $d\ge r$ and, for sufficiently small fixed $B>0$, satisfies
\begin{equation}
\label{a:eq:prop-barrier}
 \Delta b_r=Bd^{-6}(12-20\kappa r/d)
 \ge\frac{19}{2}Bd^{-6}\ge f(b_r).
\end{equation}
For example, $4\pi k\sqrt B\le19/2$ suffices.  If $R=2r$, then
\begin{equation}
\label{a:eq:prop-gap}
 b_r-b_R=B\kappa r d^{-5}\ge\gamma R^{-4}
 \quad(R\le d\le1.1R),\qquad
 \gamma=\frac{B\kappa}{2(1.1)^5}>0.
\end{equation}

\emph{Initialization.}
With probability at least $1-Cr_0^{35}$, the initial conditional density
satisfies
\begin{equation}
\label{a:eq:prop-init-good}
 P_0(y):=K*(\1_{\{d<r_0\}}\mu_0)(y)
 \le0.1\,Z/r_0\qquad(r_0\le |y|\le1.1r_0).
\end{equation}
To prove this, fix a point $y$ in this annulus and let
$A_y=\{P_0(y)>0.05Z/r_0\}$, an event of $\mathcal F_0$.
Suppose that $p=\Prob(A_y)>r_0^{40}$.  By
Lemma~\ref{lem:obs-tilt}, the raw marginal conditioned on this event
is the squared law of an antisymmetric state of offset at most
\[
 D_0+C r_0^{-2.02}\log^5(e/r_0^{40}).
\]
For the fixed $\eps$ this is at most $Z^{7/3}$ when $Z$ is large
enough.  The global estimate \eqref{eq:atomic-global-density} therefore bounds
the tilted density $\rho_{A_y}$ by
$\int\rho_{A_y}^{5/3}\le C Z^{7/3}$, with a universal $C$.

Only packet centers with $|x|<2r_0$ can contribute to the truncated
density in $P_0$: a center outside this ball has
$t_x\le c_1s^w|x|<|x|/2$, so its packet misses $B(0,r_0)$.
A truncated packet has potential at most that of the full packet,
which by radial smearing is at most $K(y-x)$.  The conditional tower
identity and H\"older's inequality consequently give
\begin{align*}
 \E[P_0(y)\mid A_y]
 &\le\int_{|x|<2r_0}\frac{\rho_{A_y}(x)}{|y-x|}\,\dd x\\
 &\le C Z^{7/5}r_0^{1/5}
   =C\eps^{6/5}\,Z/r_0.
\end{align*}
Choose $\eps$ small enough that the last coefficient is less than
$0.05$.  This contradicts the definition of $A_y$, so
$\Prob(A_y)\le r_0^{40}$.

The minimum packet width is $t_{\min}=c_1r_0^{1+w}$.  Integrating
$|y-z|^{-2}$ against a bounded packet gives the deterministic bound
\[
 \|\nabla P_0\|_\infty\le Cn t_{\min}^{-2}
 \le C Z r_0^{-2-2w}.
\]
A net in the annulus of spacing $O(r_0^2)$ has $O(r_0^{-3})$ points.
The union bound gives failure probability $O(r_0^{37})$, and the
variation between a point and its net point, divided by $Z/r_0$, is
$O(r_0^{1-2w})=o(1)$.  For large $Z$ this proves
\eqref{a:eq:prop-init-good}, with the stated weaker exponent $35$.

Let $G_0$ denote the event in \eqref{a:eq:prop-init-good}, and set
\[
 \mu^{\mathrm{in}}=\1_{G_0}\1_{\{d<r_0\}}\mu_0,
 \qquad p_0^{\mathrm{err}}=\1_{G_0^c}\1_{\{d<r_0\}}\mu_0.
\]
On $d<2r_0$ use the branch
\[
 v=V-K*\mu^{\mathrm{in}}-0.7Z/r_0
            +C_3(Z/r_0)^{3/2}d^2.
\]
Choose $C_3$ universally so that $6C_3\ge4\pi k2^{3/2}$, and then
choose $\eps$ sufficiently small to meet the preceding tilt
condition and $4C_3\eps^{3/2}\le0.05$.  The quadratic term on this
ball is at most $0.05Z/r_0$, since $Zr_0^3=\eps^3$.
For $r_0\le d<2r_0$ we thus have $v\le2Z/r_0$ and
\[
 \Delta v=-4\pi\nu+4\pi\mu^{\mathrm{in}}
                +6C_3(Z/r_0)^{3/2},
 \qquad 6C_3(Z/r_0)^{3/2}\ge f(v).
\]
On $r_0\le d<1.06r_0$, the good-event bound (and the easier bound
$\mu^{\mathrm{in}}=0$ on its complement) give
\[
 v\ge(1/1.06-0.1-0.7)Z/r_0>0.14Z/r_0.
\]
On $1.9r_0<d<2r_0$ we instead have
$v\le(1/1.9-0.7+0.05)Z/r_0<0$.
Fix $B$ satisfying the barrier condition and $0<B\le0.1\eps^3$.
Then $v>b_{r_0}$ on the lower overlap, whereas $v<0<b_{r_0}$ on
the upper overlap.  Define
\[
 u_0=\begin{cases}
 v,&d<1.06r_0,\\
 \max(v,b_{r_0}),&r_0<d<2r_0,\\
 b_{r_0},&d>1.9r_0.
 \end{cases}
\]
The definitions agree on their open overlaps.
Lemma~\ref{lem:maximum} proves \eqref{a:eq:prop-invariant} initially;
near $d=r_0$ the branch $v$ alone attains the maximum, and its displayed
Laplacian supplies both required densities.  The exterior upper bound
holds, for example, when
$C_{\mathrm{inv}}\ge\max(B,32\eps^3)$.
Also, using $n\le3Z=3\eps^3r_0^{-3}$,
\begin{equation}
\label{a:eq:prop-initial-error}
 \E\int p_0^{\mathrm{err}}
 \le n\Prob(G_0^c)\le C r_0^{32}.
\end{equation}

\emph{Order of the remaining constants.}
The constants $C_3,\eps,B$ have now been fixed universally.
Increase $C_{\mathrm{inv}}$ to exceed also the universal cap
$C_{\mathrm{cap}}$ in Proposition~\ref{a:prop:cond-comparison}.
After fixing $\gamma$ in \eqref{a:eq:prop-gap}, choose
$0<\lambda_2<\lambda_1<\gamma/2$.  Choose $L_1>2$ so large that
$C_{\mathrm{inv}}L_1^{-4}<\lambda_2$.  This ensures
\[
 C_{\mathrm{inv}}d^{-4}-\lambda_2R^{-4}<b_R
 \qquad(d>L_1R).
\]
Next fix $0<h_l<B/4$ and $h_h>C_{\mathrm{inv}}$, and then choose
$e_0,\xi>0$ sufficiently small that
\begin{equation}
\label{a:eq:prop-choices}
 2e_0(2L_1)^4<B/4,\qquad
 16\xi+2e_0<\lambda_1-\lambda_2,\qquad \xi<\lambda_2,
\end{equation}
with $\xi$ also in the range of Proposition~\ref{a:prop:cond-comparison}.
Only now restrict $s$ according to that proposition, $D_0$, and the
other smallness conditions in this proof.  Finally take $Z$ large
according to $D_0,s$ and the initialization requirements.  Thus no
constant is chosen in terms of a later scale or a particular state.

\emph{One outward step.}
Suppose that \eqref{a:eq:prop-invariant} holds at $r=r_j$, write
$R=2r$, and abbreviate $u=u_j$.  On the good band event of
Proposition~\ref{a:prop:cond-comparison}, for the band
$r\le d\le4L_1r=2L_1R$, introduce
\[
 v_1=u-\lambda_1R^{-4},\qquad
 v_2=H_j-\lambda_2R^{-4}.
\]
Define a new function on a fixed open cover by
\begin{equation}
\label{a:eq:prop-branches}
 \begin{array}{c|c}
 \text{region}&u'\\ \hline
 d<1.08R&\max(v_1,v_2)\\
 1.04R<d<2L_1R&\max(v_1,v_2,b_R)\\
 d>L_1R&b_R.
 \end{array}
\end{equation}
On bad band data omit $v_2$ in every row.  In both cases set
\begin{equation}
\label{a:eq:prop-step-error}
 \widetilde p=p_j^{\mathrm{err}}
           +\1_{\{\mathrm{bad}\}}\1_{\{r\le d<R\}}\mu_j.
\end{equation}
On the lower overlap, \eqref{a:eq:prop-gap} gives
$v_1\ge b_r-\lambda_1R^{-4}>b_R$.  On the upper overlap, the caps
on $u$ and, on good data, $H_j$ put both shifted branches below $b_R$.
The definitions therefore agree.  For $R\le d<1.08R$ the lower bound
$u'\ge b_R$ follows already from $v_1$; farther out the barrier is
included explicitly.  The upper bound $u'\le C_{\mathrm{inv}}d^{-4}$
and the equality $u'=b_R$ for $d>L_1R$ follow from the same caps.

We prove the distributional inequality at radius $R$ before averaging.
The open-neighborhood formulation is as follows.  At a point $y_0$ in
one of the open sets of \eqref{a:eq:prop-branches}, let $M$ be the maximum
of its candidate family and put $\eta=e_0R^{-4}$.  Keep precisely those
candidates $v$ with $M(y_0)-v(y_0)\le\eta$, and choose a maximizer
$v_*$ at $y_0$.  Continuity and finiteness give an open neighborhood
$U$ of $y_0$, inside the chosen open set, on which every discarded
candidate is strictly below $v_*$ and every retained candidate satisfies
\begin{equation}
\label{a:eq:prop-near-active}
 M-v<2e_0R^{-4}.
\end{equation}
The retained family has maximum $M$ throughout $U$.  At the origin
make this construction with the continuous offsets by $V$; the family
there contains only branches having that common singularity.

Each retained branch satisfies on the whole of this neighborhood the
common inequality
\begin{equation}
\label{a:eq:prop-common}
 \Delta v\ge-4\pi\nu+4\pi\1_{\{d<R\}}\mu_j
          -4\pi\widetilde p+\1_{\{d\ge R\}}f(v).
\end{equation}
For $v_1$, retain its old distributional inequality and compare its
right side with the new one almost everywhere.  On $d<r$ the old
source suffices.  On $d\ge R$, use $f(u)\ge f(v_1)$.  On the gain
region $r\le d<R$, bad data are paid by the added error in
\eqref{a:eq:prop-step-error}.  On good data, $v_2$ is a candidate there,
so \eqref{a:eq:prop-near-active} gives
\[
 H_j\le u-(\lambda_1-\lambda_2-2e_0)R^{-4}.
\]
Moreover $h_l<7B/8\le d^4u\le C_{\mathrm{inv}}<h_h$.
If $4\pi\mu_j>f(u)$, the high implication in
\eqref{a:eq:cond-inverse}, at height $h=d^4u$, would give
\[
 H_j\ge u-\xi d^{-4}\ge u-16\xi R^{-4},
\]
contrary to \eqref{a:eq:prop-choices}.  Hence
$4\pi\mu_j\le f(u)$, and the old reaction supplies the added source.

For $v_2$, which occurs only on good data, start from its identity on
all of $U$,
\[
 \Delta v_2=-4\pi\nu+4\pi\mu_j.
\]
This supplies \eqref{a:eq:prop-common} on $d<R$, since
$\widetilde p\ge0$.  At any remaining point where $v_2$ is retained,
the lower splice or the included barrier gives $M\ge b_R$, and
$d\le2L_1R$.  Thus
\[
 d^4v_2\ge B(1-\kappa R/d)-2e_0(d/R)^4
 \ge7B/8-2e_0(2L_1)^4>5B/8>h_l.
\]
The cap also gives $d^4v_2\le C_{\mathrm{cap}}<h_h$.
If $4\pi\mu_j<f(v_2)$, the low implication in
\eqref{a:eq:cond-inverse}, with $h=d^4v_2$, would imply
\[
 H_j\le v_2+\xi d^{-4}
       \le v_2+\xi R^{-4}<v_2+\lambda_2R^{-4}=H_j,
\]
which is impossible.  Consequently $4\pi\mu_j\ge f(v_2)$ there.
The barrier branch occurs only where $d>R$ and satisfies the common
inequality by \eqref{a:eq:prop-barrier}.

These verifications compare zero-order densities almost everywhere
on the same open set $U$; they do not restrict or differentiate a weak
inequality across an activity set or a sphere.  In particular they
apply on a neighborhood crossing $d=R$.  The common reaction is
$F(y,t)=\1_{\{d\ge R\}}f(t)$, which meets
Lemma~\ref{lem:maximum}'s hypotheses after canceling $V$ near
the origin.  That lemma proves \eqref{a:eq:prop-common} for the local
maximum.  These local conclusions give the inequality globally by
a partition of unity.  We have proved the desired invariant at radius
$R$ with density $\mu_j$, function $u'$, and error $\widetilde p$.

\emph{Forgetting the current observations.}
Define the next offset and error by conditional integration,
\[
 u_{j+1}=V+\E[u'-V\mid\mathcal F_{j+1}],\qquad
 p_{j+1}^{\mathrm{err}}=\E[\widetilde p\mid\mathcal F_{j+1}].
\]
We give the regularity justification as well as the formal calculation.
For the fixed system $t_{\min}>0$, and integration of each packet
gives deterministic bounds
\[
 \|\mu_j\|_\infty\le Cnt_{\min}^{-3},\qquad
 \|K*\mu_j\|_\infty\le Cnt_{\min}^{-1},\qquad
 \|\nabla K*\mu_j\|_\infty\le Cnt_{\min}^{-2}.
\]
The same potential bounds hold for the initially truncated density.
The explicit initial branches therefore have deterministic local
Lipschitz bounds for their offsets.  Subtracting a constant and taking
a finite maximum preserve such bounds.  Where $b_R$ occurs, its
offset $b_R-V$ is smooth and separated from the origin.  The fixed
open cover and agreement on overlaps preserve local bounds after
gluing.  The errors are bounded, nonnegative, and supported in a
deterministic ball.  Induction supplies all the claimed bounds at
every step, without requiring them to be uniform in $Z$.

All functions are jointly measurable in data and position.  Indeed
the conditional densities have the likelihood versions described in
Section~\ref{a:sec:conditioning}; the initialization flag is a supremum
of a continuous function on a fixed compact annulus, testable on a
countable dense set; and the band flags and finite maxima are measurable.
Conditional integration of functions with deterministic local
Lipschitz bounds preserves those bounds.  It also commutes with weak
derivatives by Fubini.  Nonnegative smooth tests in a countable dense
family first give simultaneous almost-sure weak inequalities; local
bounds extend them to every test.  Thus the following conditional
integration is legitimate as an inequality of distributions.

Integrate \eqref{a:eq:prop-common} conditionally on
$\mathcal F_{j+1}$.  The tower property supplies $\mu_{j+1}$, while
convexity gives
\[
 \E[f(u')\mid\mathcal F_{j+1}]
       \ge f\bigl(V+\E[u'-V\mid\mathcal F_{j+1}]\bigr).
\]
The deterministic bounds $b_R\le u'\le C_{\mathrm{inv}}d^{-4}$
and the equality outside $L_1R$ pass to this average.  Positivity,
boundedness, and support of the error do as well.  This proves
\eqref{a:eq:prop-invariant} at step $j+1$.

\emph{Summing the lost source.}
For $r\le d_y<2r$, every master packet contributing to $\mu_j(y)$
has its center in a fixed annulus comparable to $r$, and its width is
comparable to $t_y$.  Conditional Jensen and the annular count estimate
\eqref{a:eq:screen-annular}, applied with offset at most $D_0$, yield
\[
 \|\mu_j(y)\|_2\le Ct_y^{-3}r^{-3}
              \le Cr^{-6-3w}.
\]
Here $r\le s\le1$ and $s$ is small according to $D_0$, so
$\max(r^{-3},1)+\sqrt{D_0r}\le Cr^{-3}$.
The bad band probability is at most $Cr^{25}$ by
Proposition~\ref{a:prop:cond-comparison}.  Cauchy--Schwarz in data and
integration over the gain annulus give
\[
 \E\int(\widetilde p-p_j^{\mathrm{err}})
 \le C r^{25/2}r^3r^{-6-3w}
 =Cr^{\alpha},\qquad \alpha=\frac{19}{2}-3w>0.
\]
Conditional averaging preserves the expected integral. Stopping the
sum before any index $j$ and using \eqref{a:eq:prop-initial-error}
proves \eqref{eq:atomic-prefix-error}, with the same master kernels
and without replacing an expectation by a pathwise bound. At the
terminal index the geometric progression of radii gives
\begin{equation}
\label{a:eq:prop-total-error}
 \int p_m^{\mathrm{err}}
 =\E\int p_m^{\mathrm{err}}
 \le Cr_0^{32}+C\sum_{j=0}^{m-1}r_j^\alpha
 \le C(s^{32}+s^\alpha).
\end{equation}
The last error is deterministic because $\mathcal F_m$ is trivial.
The constants are independent of the number of scales and of $Z$.

\emph{From the potential to the raw electron count.}
Put $r=r_m$, $u=u_m$, and introduce the compactly supported signed
measure
\[
 S=\nu-\1_{\{d<r\}}\mu_m\,\dd y+p_m^{\mathrm{err}}\,\dd y.
\]
The invariant implies $\Delta(u-K*S)\ge0$.  Its singularities cancel
in the identity
\[
 u-K*S=(u-V)+K*(\1_{\{d<r\}}\mu_m)-K*p_m^{\mathrm{err}},
\]
whose right side is continuous and locally $H^1$ on all of $\R^3$.
The difference tends to zero at infinity: $u=b_r$ outside $L_1r$,
and the potential of the finite compact source $S$ tends to zero.
For each $a>0$, the positive part $(u-K*S-a)_+$ has compact support
and is an admissible $H^1$ test in its subharmonic inequality.  Testing
gives a nonpositive squared gradient integral, hence this positive
part vanishes.  Letting $a\downarrow0$ proves $u\le K*S$.

On the sphere $d=2r$, the invariant gives
$u\ge b_r(2r)=15B r^{-4}/256$.  The spherical mean of $K*S$
there equals the total mass of $S$ divided by $2r$, since $S$ is
supported in $\{d\le r\}$.  Consequently
\begin{equation}
\label{a:eq:prop-charge}
 Z-\int_{d<r_m}\mu_m+\int p_m^{\mathrm{err}}
 \ge\frac{15B}{128}r_m^{-3}.
\end{equation}
No sign condition on the measure $S$ is needed for this mean identity.

Finally $\mu_m$ is the unconditional master-smoothed density.
If $|x|<s/4$, then $r_0\le s/2$ and
$D_x^0\le s/2$, so $t_x\le c_1(s/2)s^w<s/4$.
The packet at $x$ is therefore wholly contained in
$B(0,s/2)\subset B(0,r_m)$.  Since every packet has mass one,
\[
 \E_\Psi\#\{i:|x_i|<s/4\}\le\int_{d<r_m}\mu_m.
\]
Combining this with \eqref{a:eq:prop-charge}, \eqref{a:eq:prop-total-error},
and $r_m\le s$ gives
\[
 Z-\E_\Psi\#\{i:|x_i|<s/4\}
 \ge\frac{15B}{128}s^{-3}-C(s^{32}+s^\alpha).
\]
Take $c_{\mathrm{def}}=15B/256$ and decrease $s_*(D_0)$ so that
the error is at most $c_{\mathrm{def}}s^{-3}$ and
$c_{\mathrm{def}}s^{-3}>1$ for every $0<s\le s_*(D_0)$.
All remaining requirements are met by increasing $Z_*(D_0,s)$.
This proves the proposition.
\end{proof}
